\documentclass[11pt]{article}
\usepackage[T1]{fontenc}
\usepackage{lmodern}
\usepackage[margin=0.85in]{geometry}
\usepackage{amsmath,amssymb,amsthm}
\usepackage{booktabs}
\usepackage{tabularx}
\usepackage{array}
\usepackage{microtype}
\makeatletter
\renewcommand\@listI{\leftmargin\leftmargini\parsep 0pt\topsep 3pt plus 1pt\itemsep 1pt plus 1pt}
\let\@listi\@listI
\makeatother
\usepackage[hidelinks]{hyperref}
\hypersetup{pdftitle={Quantum Relational Causality: Rank-Three Selection from a Nested-Reciprocity Source Law, and Conditional Gauge Structure},pdfauthor={Hendrixx}}

\newtheorem{theorem}{Theorem}[section]
\newtheorem{proposition}[theorem]{Proposition}
\newtheorem{lemma}[theorem]{Lemma}
\newtheorem{corollary}[theorem]{Corollary}
\theoremstyle{definition}
\newtheorem{definition}[theorem]{Definition}
\theoremstyle{remark}
\newtheorem{remark}[theorem]{Remark}

\newcommand{\Spre}{S^{\mathrm{NR}}_{\mathrm{pre}}}
\newcommand{\Z}{\mathbb{Z}}
\newcommand{\R}{\mathbb{R}}
\newcommand{\C}{\mathbb{C}}
\newcommand{\ch}{\operatorname{ch}}
\newcommand{\Aut}{\operatorname{Aut}}
\newcommand{\doc}[1]{\nolinkurl{#1}}
\makeatletter\g@addto@macro\UrlBreaks{\do\_}\makeatother
\newcolumntype{Y}{>{\raggedright\arraybackslash}X}

\title{Quantum Relational Causality: Rank-Three Selection from a\\
Nested-Reciprocity Source Law, and Conditional Gauge Structure}
\author{Hendrixx (independent researcher)}
\date{Version 1.0.0, November 2026}

\begin{document}
\maketitle

\begin{abstract}
Quantum Relational Causality (QRC) investigates discrete relational actions on matching
groupoids, treating spatial distance as the inferred measurement of causal impedance
between relational excitations rather than an absolute background container. This preprint
states what the program has established mathematically, under which explicit assumptions,
and what remains open. The main result concerns a declared source law $\Spre$ on a groupoid
of finite sets $D$ of excitation types with transverse matchings between $D$ and its binary
spans. For activity $0<h<2\lambda$ the unique ground state up to relabelling has $|D|=3$
and the complement matching, with excitation gap $\min(3h,2\lambda-h)$. Two further declared
laws (least endpoint completion and a centred unit-rate response) then give the endpoint group
$\Z[D]$ of rank three, exact word-ball counts $V_3(R)=(4R^3+6R^2+8R+3)/3$, return
probability $e^{-t}I_0(t/3)^3$, spectral dimension three and one co-metric
$q^\sharp=\tfrac13\sum_{d\in D}d\otimes d$. These are implication theorems for a declared
combinatorial law; the law is not derived from informal ontological principles, and nothing
here shows that nature obeys it. Separately, on a supplied Hermitian $2+3$ internal carrier
with one determinant record, the stabilizer is $(SU(2)\times SU(3)\times U(1))/\Z_6$, one
anomaly-free 15-state chiral family follows, and a pair-dual anomaly polynomial
$(m+n-5)\ch_3(E)$ selects the ranks $2+3$ within its own hypotheses; dynamical selection
of the carrier, mirror removal and a joint gauge continuum are not proved. We record
corrections to earlier claims, including six that failed independent audit, and fit the
remaining cosmology candidates to public supernova, BAO and chronometer data: one
phenomenological causal-memory dark-energy model fits at the level of $\Lambda$CDM, and
the data neither prefer nor exclude it. Every numerical statement is reproducible with
scripts in the accompanying package, each shown to fail on deliberately wrong input.
\end{abstract}

%======================================================================
\section{Introduction and scope}
\label{sec:intro}

The conceptual motivation behind Quantum Relational Causality (QRC) is relational:
seeking to understand whether spatial dimension, gauge symmetries, and metric structure
can emerge from discrete relational configurations without assuming a prior continuous
background manifold. Informally, the program explores whether the experienced
three-dimensional spatial world can be understood as an assembled hierarchy of relational
excitation grades (from a zero-dimensional base condition to primitive lines and spans),
with time as the history of causal updates.

Crucially, this preprint does not claim to deduce physical reality from that informal
metaphor, nor are the mathematical theorems derived from it. Instead, the program
organises its technical proof obligations into a sequence of \emph{gates}
(Section~\ref{sec:gates}), where each gate defines an explicit mathematical hypothesis and
a verifiable pass condition.

This preprint has a narrow purpose: to state exactly which of those obligations have
been met as mathematics, under which hypotheses, and which have not. It is written for
readers who want to check the work, not for readers who want to be persuaded by it.

\paragraph{What is claimed.}
\begin{enumerate}
\item For the declared source law $\Spre$ (Definition~\ref{def:action}), rank three with
  the complement matching is the unique ground state in the open window $0<h<2\lambda$,
  with an exact gap (Theorem~\ref{thm:selection}). This is proved here in full and checked
  by exhaustive enumeration (Section~\ref{sec:g1numerics}).
\item Under two further declared laws (least endpoint completion and a centred response),
  the selected carrier has cubic word-ball growth, the stated heat kernel, spectral
  dimension three and a single infrared co-metric (Sections~\ref{sec:endpoint}
  to~\ref{sec:cometric}).
\item On a supplied internal carrier, the stated gauge group, family decomposition,
  anomaly cancellation and pair-dual rank selection (Section~\ref{sec:gate5}).
\end{enumerate}

\paragraph{What is not claimed.} That the source law follows from the foundational
statement or from nonlocality; that any QRC law is realised in nature; that word length
or diffusion time has been identified with measured distance or elapsed time; Lorentzian
spacetime; a quantum continuum; gravity; gauge dynamics; a particle spectrum; an absolute
length scale. Section~\ref{sec:open} lists these as open problems.

\paragraph{Claim discipline.} A theorem proved after adopting an explicit source or
quantisation law is an implication theorem for that law. An implication theorem for a
declared law is not evidence that nature obeys the law, and it is not evidence that more
generic premises force the law. Where this paper writes ``Gate~1 is passed'' it means
``passed as mathematics for the declared source law'', which is the wording of the
project status record.

\paragraph{Evidence labels.} ``Proved'' means a complete proof is given in this paper.
``Sketched'' means the argument is outlined here and written out in a repository
document named by file in Appendix~\ref{app:repro}; the author's repository is not yet
public, so sketched results should be read as claims until those documents are released.
``Checked'' means executed by a package script that is also shown failing on
deliberately wrong input. ``Cited'' means taken from the literature, not re-derived.

%======================================================================
\section{Conceptual motivation and operational definitions}
\label{sec:foundation}

The informal conceptual heuristic that originally motivated the project is stated as:

\begin{quote}
Space in its underlying condition is 0D and nonlocal; particle fields are 1D
excitations; electromagnetic excitations are 2D; their composition supports the
experienced 3D world, retaining its lower-dimensional constituents; and time is the
fourth dimension encompassing its history. The limiting causal-update and propagation
role is assigned to $c$.
\end{quote}

The governing operational measurement definition is stated exactly:

\begin{quote}
Distance is the inferred measurement of causal impedance. Causal impedance does not
cause a separate quantity or substance called distance.
\end{quote}

\begin{remark}[Non-axiomatic status]
The dimensional narrative above is an informal conceptual heuristic, not a mathematical
premise. Nothing in this paper is mathematically deduced from the informal phrasing that
particle fields are ``1D'' or that electromagnetism is ``2D''. In standard mathematical
physics, physical fields are operator distributions or sections of fiber bundles over
spacetime, not lower-dimensional objects composing into space. Within the formal
mathematics of Section~\ref{sec:gate1}, the numerals $0,1,2,3$ attached to relational grades
denote filtration depths in an exterior matching complex, while the spatial dimension
three of the conclusion is the operational exponent of a declared walk and ball count.
\end{remark}

These are starting perspectives. Their mathematical realisation and physical adequacy are what
the gates must establish. ``Nonlocal'' here means that spatial locality has not yet been
inferred, not that a superluminal speed has been measured. No background metric,
elementary graviton or universal synchronised update clock is supplied by the statement.

\paragraph{Response data.} Let $H$ denote a relational configuration or history and
$Z_H$ its causal-impedance response data. The type of $Z_H$ is deliberately not fixed in
advance: it is not assumed to be an electrical resistance, a scalar edge length, an
update frequency or a dissipative coefficient. Gate~0 must specify its mathematical type
and composition law.

\paragraph{Measurement definition.} For a declared measurement protocol $P$,
\begin{equation}
d_P(a,b) := \mathcal M_P[Z_H; a,b].
\label{eq:measurement}
\end{equation}
This is a measurement definition, not an equation of motion, and no particular map
$\mathcal M_P$ has been derived. It is not permitted to choose $Z_H$ from an independently
supplied Euclidean distance and then present the recovered distance as a derivation.

\paragraph{Three meanings of ``dimension''.} The formulation keeps apart (i) the
dimensionality of a primitive excitation, (ii) the algebraic or form degree used to
describe it and (iii) the operational spatial or spectral dimension of the assembled
world; form degree alone proves neither an ontology nor the dimension of space. In
Section~\ref{sec:gate1}
the numerals $0,1,2,3$ attached to relational grades denote composition depth, and the
``three'' of the conclusion is an operational exponent of a declared walk and ball count.
Identifying that exponent with the dimension of physical space requires the response and
clock bridge of Gates~1 and~2, which is open.

%======================================================================
\section{The gate structure and current status}
\label{sec:gates}

The gates are proof dependencies, not stages that the universe passed through in an
external time.

Gate~0 specifies the nonlocal 0D condition (states, excitations, response observables,
composition rules) with causal impedance defined without a prior distance. Gate~1
recovers measured three-dimensional space: cubic volume growth and spectral dimension
three, with metric, measure, operator and clock identified as outputs of the declared
response. Gate~2 recovers time, the causal limit $c$ and one Lorentzian cone shared by
every physical mode. Gate~3 recovers quantum preparation, evolution and measurement,
including $p(i)=\operatorname{Tr}(\rho E_i)$ and the Tsirelson bound. Gate~4 recovers
gravity as the matter-dependent response of causal impedance, classically and in a
quantum completion. Gate~5 identifies the electromagnetic, weak and strong sectors in one
interacting phase. Gate~6 derives the particle spectrum and discriminating predictions.
A seventh stage collects cosmological and other observational consequences.

Table~\ref{tab:status} reproduces the status table of the current integrated position
(project status record, review completed 2026-11-19). Dated sections of that record
written before this position, in particular the September 2026 entries headed
``Level-4'' to ``Level-7'' (triality generations, spectral triples, asymptotic safety,
horizon-impedance dark energy, leptogenesis and others), are historical. None of them is
claimed here; Section~\ref{sec:corrections} reports audits of several.

\begin{table}[ht]
\centering\small
\begin{tabularx}{\textwidth}{@{}c Y Y@{}}
\toprule
Stage & What the reviewed mathematics supports & Open physical dependency \\
\midrule
0 & Relabelling-invariant premetric interface and candidate dynamics & Selection of a
source and physical history law \\
1 & Exact rank selection under transverse reciprocity, with cubic growth and diffusion
under the specified completion and walk & Source selection and operational
identification of the clock and ruler \\
2 & Exact bracket tests for the declared scalar coefficients, finite cochain
certificates and conditional continuum results & One complete interacting mode inventory
and a common physical cone \\
3 & Finite record/Born construction, separated marginal tests and CHSH bound &
Operational frequency law on the same positive interacting histories \\
4 & Conditional ADM coefficient and traceless kinetic-form calculation & Positive
two-helicity interacting response and universal stress coupling \\
5 & Finite exterior family, quotient descent, anomaly cancellation and charge
decomposition on a supplied internal seed & Dynamical selection, mirror removal and a
joint gauge continuum \\
6 & Allowed mass channels, exact one-family seesaw and freedom of flavour matrices &
Generation count, matrix entries, vacuum and physical spectrum \\
7 & Conditional observable calculations and a constant-equation-of-state bounce for the
displayed effective law & Source-derived predictions, independent likelihoods and
continuum-to-poset return map \\
\bottomrule
\end{tabularx}
\caption{Gate status from the project status record (current integrated position,
2026-11-19), reproduced verbatim. ``Certificate'' there means a finite computation on a
declared object, not an external certification.}
\label{tab:status}
\end{table}

The status record summarises Gates~1 and~5 with these tags:
\begin{center}\small\ttfamily
\begin{tabular}{@{}l@{}}
MATHEMATICAL\_GATE1\_PASSED\_FOR\_DECLARED\_SOURCE\_LAW\\
EMPIRICAL\_SOURCE\_LAW\_UNCONFIRMED\\
ABSOLUTE\_SCALE\_UNDETERMINED\\
CONDITIONAL\_INTERNAL\_GROUP\_AND\_ANOMALY\_THEOREMS\_AVAILABLE\\
FULL\_GATE5\_NOT\_PASSED
\end{tabular}
\end{center}
Gates 2, 4, 5 and 6 are marked not passed. Gate~3 is marked closed only within an
adopted sharp-record quantum law.

%======================================================================
\section{Gate 1: rank selection and the cubic-diffusive phase}
\label{sec:gate1}

\subsection{State space and source law}
\label{sec:state}

\begin{definition}[Source groupoid]
\label{def:groupoid}
An object is a pair $(D,M)$ where
\begin{enumerate}
\item $D$ is a finite unlabelled set of primitive reversible excitation types;
\item $P_2(D)=\binom{D}{2}$ is the set of all unordered binary spans (every pair is present);
\item $M\subseteq D\times P_2(D)$ is a partial matching, injective in both projections;
\item every matched pair is transverse: $(d,p)\in M \Rightarrow d\notin p$.
\end{enumerate}
Morphisms are bijections of $D$ that carry $M$ to $M$. The empty object $D=\varnothing$
represents the undifferentiated 0D condition: it has no move, adjacency or distance.
\end{definition}

No coordinate, metric, orientation, rank or length is part of an object. The relational
grades $\Lambda^0,\dots,\Lambda^3$ of $\R[D]$ (undifferentiated condition, primitive
excitation, coupled binary span, transverse interleaving) form an increasing filtration,
so every higher grade retains the lower ones; the numerals are composition depth.

\begin{definition}[Nested-reciprocity activity law]
\label{def:action}
Write $n=|D|$ and $m=|M|$. The reciprocity defect is
$\Delta(D,M)=(n-m)+\bigl(\binom n2-m\bigr)$, the number of unmatched primitives plus the
number of unmatched spans. For $\lambda>0$ and $h\in\R$,
\begin{equation}
\Spre(D,M)=\lambda\,\Delta(D,M)-h\,n .
\label{eq:action}
\end{equation}
The groupoid partition function is
$Z_\beta=\sum_{n\ge0}\frac{1}{n!}\sum_{M}e^{-\beta \Spre([n],M)}$, where the inner sum runs
over labelled matchings on $[n]$; the factor $1/n!$ gives each isomorphism class the
weight $1/|\Aut(D,M)|$.
\end{definition}

This is schema \texttt{qrc.spre.nested\_reciprocity\_activity.v1} in the repository. The
load-bearing physical assumptions of the selection result are:
\begin{enumerate}
\item[(H1)] every unordered binary span of excitation types is present;
\item[(H2)] matching is transverse and one-use (injective on both sides);
\item[(H3)] unmatched primitives and unmatched spans both cost a positive amount (equal
  costs in~\eqref{eq:action}; Proposition~\ref{prop:unequal} relaxes this);
\item[(H4)] the activity lies in the open window $0<h<2\lambda$.
\end{enumerate}
None of (H1) to (H4) is derived from the foundational statement. The repository records
explicitly that generic nonlocality does not force this law.

\subsection{The selection theorem}
\label{sec:selection}

\begin{lemma}[Maximal transverse matching]
\label{lem:maxmatch}
The largest transverse matching has size $0$ for $n\le2$ and size $n$ for every $n\ge3$.
Hence $\min_M\Delta=0,1,3$ for $n=0,1,2$ and $\min_M\Delta=\binom n2-n=n(n-3)/2$ for
$n\ge3$.
\end{lemma}

\begin{proof}
For $n=1$ there are no spans; for $n=2$ the only span contains both elements, so
transversality forbids every match. For $n\ge3$ place $D$ temporarily in a cyclic order
$d_1,\dots,d_n$ and match $d_i\mapsto\{d_{i+1},d_{i+2}\}$ (indices mod $n$). The target
spans are distinct and omit their sources, so this is a transverse matching of size $n$.
No matching is larger because the projection to $D$ is injective. The cyclic order is
discarded afterwards. Substituting $m$ into $\Delta$ gives the minima.
\end{proof}

In particular $\Delta=0$ exactly when $n=0$, or $n=3$ and $M=\{(d,D\setminus\{d\}):d\in D\}$:
at $n=3$ each element is transverse only to its complementary pair, so the lossless
matching is forced.

\begin{theorem}[Rank selection and gap]
\label{thm:selection}
The ground states of~\eqref{eq:action} over all objects of all ranks are: the empty
object alone if $h<0$; the empty object and the rank-three complement matching if $h=0$;
the rank-three complement matching alone, up to relabelling, if $0<h<2\lambda$. In the
last case the gap to the first excited object is exactly
\begin{equation}
g=\min\{3h,\;2\lambda-h\}>0 .
\label{eq:gap}
\end{equation}
\end{theorem}

\begin{proof}
Since $\Spre$ decreases in $m$ at fixed $n$, the minimum at rank $n\ge3$ is
$E_n=\frac{\lambda}{2}n(n-3)-hn$ by Lemma~\ref{lem:maxmatch}. Then $E_3=-3h$ and, for
$n\ge4$,
\begin{equation}
E_n-E_3=(n-3)\Bigl(\frac{\lambda n}{2}-h\Bigr),
\label{eq:diff}
\end{equation}
which on $0<h<2\lambda$ is positive and smallest at $n=4$, where it equals $2\lambda-h$.
The empty object has energy $0=E_3+3h$. Rank one has energy $\lambda-h=E_3+\lambda+2h$,
rank two $3\lambda-2h=E_3+3\lambda+h$, and removing one match from the rank-three ground
state costs $2\lambda$. Each of these is at least $\min\{3h,2\lambda-h\}$: if
$h\le\lambda/3$ then $3h\le\lambda+2h$, and if $h\ge\lambda/3$ then
$2\lambda-h\le\lambda+2h$. For $h<0$ every nonempty object has $\Spre\ge-hn>0$, and at
$h=0$ the zero-energy objects are those with $\Delta=0$.
\end{proof}

\begin{proposition}[Unequal costs]
\label{prop:unequal}
Replace~\eqref{eq:action} by $\lambda_1(n-m)+\lambda_2\bigl(\binom n2-m\bigr)-hn$ with
$\lambda_1,\lambda_2>0$. For $0<h<2\lambda_2$ the rank-three complement matching is the
unique ground state and the gap is
$g_{12}=\min\{3h,\ \lambda_1+2h,\ \lambda_1+\lambda_2,\ 2\lambda_2-h\}$. The window is set
by $\lambda_2$, the cost of unmatched spans.
\end{proposition}

\begin{proof}[Proof (complete)]
At maximal matching $E_n=\frac{\lambda_2}{2}n(n-3)-hn$ for $n\ge3$, so $E_{n+1}-E_n=
\lambda_2(n-1)-h>0$ for $n\ge3$ on the window. The empty, rank-one, rank-two,
one-missing-match and rank-four gaps are $3h$, $\lambda_1+2h$, $2\lambda_1+\lambda_2+h$,
$\lambda_1+\lambda_2$ and $2\lambda_2-h$; the rank-two value exceeds the
one-missing-match value, and every further missing match adds $\lambda_1+\lambda_2$.
\end{proof}

\begin{proposition}[The activation boundary does not select activation]
\label{prop:boundary}
At $h=0$ the zero-temperature groupoid probabilities are $\Pr(n=3)=1/7$ and
$\Pr(n=0)=6/7$.
\end{proposition}

\begin{proof}
The two zero-energy classes have $|\Aut|=1$ (empty) and $|\Aut|=|S_3|=6$ (complement
matching, which every permutation preserves). Their weights are $1$ and $1/6$.
\end{proof}

So neither zero defect nor groupoid entropy derives the activated phase; $h>0$ is an
assumption. (A repository variant inserts one abstract excitation as a boundary source;
that lifts the degeneracy but explains neither the source nor a transition out of the
empty sector.)

\begin{theorem}[All-rank Gibbs concentration; sketched]
\label{thm:gibbs}
$Z_\beta<\infty$ for every $\beta>0$, and for each $\beta_0>0$ there is a finite
$C_{\beta_0}$ with $\Pr_\beta[(D,M)\not\cong\text{ground}]\le C_{\beta_0}e^{-(\beta-\beta_0)g}$
for $\beta\ge\beta_0$.
\end{theorem}

\begin{proof}[Sketch]
The number of labelled partial matchings at rank $n$ is at most
$\sum_{m\le n}\binom nm\binom{\binom n2}{m}m!=e^{O(n\log n)}$, while the minimal action grows
as $\lambda n^2/2$. The sum therefore converges, and every excited term carries an extra
factor $e^{-(\beta-\beta_0)g}$ relative to its $\beta_0$ value.
\end{proof}

\begin{remark}[Where the three comes from]
\label{rem:three}
No numeral three is placed in~\eqref{eq:action}; it comes from (H1), since lossless
reciprocity requires $n=\binom n2$. With $k$-ary spans the condition is $n=\binom nk$,
solved by $n=k+1$ (the check confirms that ternary spans select rank four). The result
reads ``binary coupling plus lossless transverse reciprocity selects three excitation
types'' and stands or falls with that premise. Three types are not yet a spatial
dimension; that needs Sections~\ref{sec:endpoint} to~\ref{sec:heat}.
\end{remark}

\subsection{Least endpoint completion}
\label{sec:endpoint}

\begin{definition}[Endpoint law]
Signed finite words in $D$ are identified under the least congruence generated by
$dd^{-1}\sim1$ and $de\sim ed$ for $d,e\in D$ (each interchange is witnessed by the span
$\{d,e\}$). ``Least'' means that no other identification is imposed.
\end{definition}

The resulting endpoint group is the universal abelian group on $D$: $A(D)=\Z[D]$, free
abelian of rank $|D|$. Equivalently, rotor characters $D\to SO(2)$ separate distinct
exponent vectors. Relational positions form an $A(D)$-torsor, so no origin is chosen.
The word distance is $r_D(x,y)=\sum_d|k_d|$ for $y-x=\sum_dk_d[d]$; it is invariant
under translation, reversal of each primitive and every relabelling. Finite regulators
are the characteristic quotients $Q_m=A(D)/mA(D)$, which choose no axes and have word
systole exactly $m$.

The completion clause is load bearing, and each alternative is a countermodel:
the extra relation $\sum_{d}d=0$ reduces the rank to two and gives quadratic growth;
omitting interchange gives the free non-abelian group and exponential growth; letting the
three types act trivially gives no extended space at all.

\subsection{Word-ball growth}
\label{sec:volume}

\begin{proposition}[Exact ball count]
\label{prop:volume}
For $|D|=n$, the number of elements of $A(D)$ within word radius $R$ is
$V_n(R)=\sum_{j=0}^{n}2^j\binom nj\binom Rj$. In the selected phase
\begin{equation}
V_3(R)=\frac{4R^3+6R^2+8R+3}{3}\;\sim\;\frac43R^3 ,
\label{eq:v3}
\end{equation}
and on $Q_m$ the same count holds exactly for $2R<m$.
\end{proposition}

\begin{proof}
Choose the $j$ nonzero coefficients ($\binom nj$ ways), their signs ($2^j$) and positive
magnitudes with sum at most $R$ ($\binom Rj$ ways). Summing for $n=3$ gives
\eqref{eq:v3}. On $Q_m$ the quotient map is injective on balls of radius $R$ when the
systole $m$ exceeds $2R$.
\end{proof}

Equation~\eqref{eq:v3} is the standard lattice-point count of the octahedron
$|x|+|y|+|z|\le R$ in $\Z^3$; the mathematics is classical, and the content of
Gate~1 lies in the chain that leads to it. Two limits on what it shows: no bounded-radius
test can certify asymptotic dimension (the map $\Z^4\to\Z^3$,
$e_4\mapsto\sum_iN_ie_i$ with $1+\sum_i|N_i|>2R$, is injective on every radius-$R$ ball),
and a constant physical length per step, if one is later assigned, changes only the
coefficient $4/3$.

\subsection{Centred response, heat kernel and spectral dimension}
\label{sec:heat}

Declare the response generator to be (R1) conservative and translation invariant,
(R2) invariant under reversal of each primitive and under relabelling, (R3) supported on
single primitive moves $\pm[d]$, (R4) free of killing, drift and long jumps, and (R5)
normalised to total jump rate one. The empty object has generator $0$ and return
probability $1$.

\begin{proposition}[Unique generator]
For nonempty $D$, (R1) to (R5) give exactly
\begin{equation}
L_Df(x)=\frac{1}{2|D|}\sum_{d\in D}\bigl[f(x+[d])+f(x-[d])-2f(x)\bigr].
\label{eq:generator}
\end{equation}
\end{proposition}

\begin{proof}
By (R1) and (R3) the generator is determined by rates $r_{\pm d}$. Reversal invariance
gives $r_{+d}=r_{-d}$, relabelling invariance makes all rates equal, and (R5) fixes the
common value $1/(2|D|)$.
\end{proof}

(R1) to (R5) are a declared clock class; the parameter $t$ below is a dimensionless
diffusion parameter, not physical time (Gate~2).

\begin{proposition}[Return probability]
On the infinite cover,
\begin{equation}
p_D(t;0,0)=e^{-t}\,I_0\!\left(t/|D|\right)^{|D|}.
\label{eq:heat}
\end{equation}
\end{proposition}

\begin{proof}
Under~\eqref{eq:generator} the $|D|$ coefficients evolve as independent continuous-time
symmetric walks on $\Z$, each with total jump rate $1/|D|$. A symmetric walk on $\Z$ with
total rate $\nu$ returns with probability $e^{-\nu t}I_0(\nu t)$; take the product.
\end{proof}

\begin{theorem}[Spectral dimension three]
\label{thm:ds}
For $|D|=3$,
\begin{equation}
d_s(t):=-2t\frac{d}{dt}\log p_D(t;0,0)=2t\Bigl[1-\frac{I_1(t/3)}{I_0(t/3)}\Bigr]
\;\longrightarrow\;3\qquad(t\to\infty).
\label{eq:ds}
\end{equation}
\end{theorem}

\begin{proof}
Differentiate~\eqref{eq:heat} using $I_0'=I_1$. The large-argument expansion
$1-I_1(x)/I_0(x)=\frac{1}{2x}+O(x^{-2})$ with $x=t/3$ gives $d_s(t)=3+O(1/t)$. The
logarithmic definition $-2\log p/\log t$ also tends to $3$, since
$p_D\sim(3/2\pi t)^{3/2}$.
\end{proof}

\paragraph{Finite regulators.} On $Q_m$ a return that is not a return on the cover needs
a displacement in $mA(D)\setminus\{0\}$, hence at least $m$ jumps. So
$|p_m(t)-p_\infty(t)|\le\Pr\{\mathrm{Pois}(t)\ge m\}$, and termwise differentiation gives
$|p_m'(t)-p_\infty'(t)|\le2\Pr\{\mathrm{Pois}(t)\ge m-1\}$. Along $t_m=\sqrt m$ both are
superpolynomially small compared with $t_m^{-3/2}$, so $d_{s,m}(t_m)\to3$. A fixed finite
quotient eventually mixes and has late-time spectral dimension $0$; the cofinal order of
limits is part of the statement.

\paragraph{Joint source and heat limit (sketched).} Instead of selecting the ground state
first, take $t_m=\sqrt m$ and $\beta_m=\beta_0+\frac2g\log m$. Theorem~\ref{thm:gibbs}
gives $\Pr_{\beta_m}(\text{not ground})=O(m^{-2})=O(t_m^{-4})$. Every nonempty normalised
generator has return in $[0,1]$ and a bounded time derivative, and the empty generator
has return $1$ and derivative $0$, so the annealed derivative spectral dimension, taken
at fixed $(\beta,m)$ and then evaluated on the diagonal, also tends to $3$. The integer
window $R_m=\lfloor\sqrt m\rfloor$ stays prewrap, so~\eqref{eq:v3} holds there with
probability tending to one.

\subsection{The infrared co-metric}
\label{sec:cometric}

\begin{proposition}[Acoustic co-metric]
\label{prop:cometric}
For a character $\theta\in\operatorname{Hom}(A(D),\R)$ the Bloch eigenvalue of $-L_D$ at
$|D|=3$ is
\begin{equation}
\varepsilon_D(\theta)=\frac13\sum_{d\in D}\bigl(1-\cos\theta(d)\bigr)
=\frac12\,q^\sharp(\theta,\theta)+O(|\theta|^4),\qquad
q^\sharp=\frac13\sum_{d\in D}d\otimes d .
\label{eq:cometric}
\end{equation}
\end{proposition}

\begin{proof}
Apply~\eqref{eq:generator} to $e^{i\theta(x)}$ and expand $1-\cos u=u^2/2+O(u^4)$.
\end{proof}

The tensor $q^\sharp$ is coordinate free and invariant under relabelling and
reversal. Three elementary consequences: (i) in the one-vertex periodic cell each
unoriented edge pair has conductance $1/3$ and the corrector incidence vanishes, so the
effective tensor is $A_{\mathrm{eff}}=q^\sharp$, and $K_q=-2L_D$ has principal symbol
$q^\sharp$; this fixes the factor of two between the Markov clock and the wave
normalisation, so word distance, volume, diffusion and the scalar seed use one metric.
(ii) The $q$-ball count is $\#\{a: q(a,a)\le R^2\}=\frac{4\pi}{9\sqrt3}R^3+O(R^2)$, with
$q=(q^\sharp)^{-1}$. (iii) A scalar Hamiltonian built on $c^2K_q$ has infrared dispersion
$\omega^2=c^2q^\sharp(k,k)+m^2+O(|k|^4)$. Item (iii) is only a principal seed for Gate~2:
it does not show that every physical species shares this symbol, and $q^\sharp$ selects
an $O(3)$ frame orbit, not an oriented coframe or a Lorentzian history.

\subsection{Numerical verification}
\label{sec:g1numerics}

Two scripts check the Gate~1 statements. Each records an expected outcome for every
test, including tests on deliberately wrong inputs that must fail, and exits with status
$0$ only if every observed outcome matches its expectation. Table~\ref{tab:g1checks}
summarises them; full output is in the package.

\begin{table}[ht]
\centering\small
\begin{tabularx}{\textwidth}{@{}Y Y Y@{}}
\toprule
Statement & Independent computation & Result \\
\midrule
Theorem~\ref{thm:selection}, Proposition~\ref{prop:unequal} &
Exhaustive enumeration of every labelled object for $n\le6$ (772{,}621 objects at
$n=6$), exact rational arithmetic, 11 parameter points; bound $m\le n$ for
$6<n\le400$ & Unique ground state is the rank-three complement matching, $|\Aut|=6$;
gaps equal~\eqref{eq:gap} and $g_{12}$ exactly \\
Lemma~\ref{lem:maxmatch} & Enumeration ($n\le6$), augmenting-path matching ($n\le12$) &
max $m=0,0,0,3,4,\dots,12$ \\
Enumerator & Independent bitmask dynamic programme & counts agree for $n\le6$ \\
Proposition~\ref{prop:boundary} & Groupoid weights of the zero-energy classes &
$\Pr(n=3)=1/7$ exactly \\
Equation~\eqref{eq:v3} & Breadth-first search on the Cayley graph of $\Z[D]$, $R\le30$;
torus $Q_{15}$ & exact; fails after wrap ($R>7$) as stated \\
Equation~\eqref{eq:heat} & Poissonised multinomial series; $\exp(tL_D)$ on a $41^3$
torus & relative error $3\times10^{-15}$ and $8\times10^{-16}$ \\
Theorem~\ref{thm:ds} & Finite-difference derivative of $\log p$ & $d_s=3.362$, $3.023$,
$3.002$, $3.0002$, $3.00002$ at $t=10$ to $10^5$ \\
Equation~\eqref{eq:cometric} & Symbol of $L_D$ on plane waves, random non-orthogonal
realisation of $D$ & relative error $1.6\times10^{-7}$ \\
\bottomrule
\end{tabularx}
\caption{Gate~1 checks. Rows 1 to 4: \texttt{check\_gate1\_selection.py} (47 tests, 14
of them on wrong inputs). Rows 5 to 8: \texttt{check\_gate\_results.py} (from an
independent audit; 26 tests, 13 on wrong inputs, covering Gate~5 as well).}
\label{tab:g1checks}
\end{table}

The wrong inputs, all of which fail as required, are: $h=2\lambda$, $h>2\lambda$, $h=0$
and $h<0$; dropping transversality (six labelled lossless matchings at rank three);
ternary spans (rank four selected); no cost for unmatched spans (no minimum); a window
set by $\lambda_1$ instead of $\lambda_2$; wrong gap formulas; labelled instead of
groupoid weights; and an enumerator without span-side injectivity.

\subsection{What Gate 1 does not establish}
\label{sec:g1scope}

\begin{enumerate}
\item The source law (H1) to (H4), the endpoint law and the response class (R1) to (R5)
  are declared, not derived. The repository states that generic nonlocality does not
  force them.
\item The selected ``three'' is the rank of an endpoint group and the exponent of a
  declared ball count and walk. Identifying word length with measured distance and the
  diffusion parameter with elapsed time requires a protocol $\mathcal M_P$ in
  \eqref{eq:measurement} and a clock, which are open (Section~\ref{sec:open}).
\item Positive activity is assumed (Proposition~\ref{prop:boundary}).
\item No absolute length is determined, and none is needed for the exponents.
\item The repository's complete Gate~1 contract (schema version 2, ``refinement
  stationary'') replaces the characteristic quotients by random prime-cyclic regulators
  with a Haar-distributed quenched response. Its additional premises (strong monoidality
  and split-and-recombine stationarity of the response, quenched response, independent
  additive gluing of singleton attempt clocks) are load bearing, are not derived from the
  rank action, and are not reviewed here. The characteristic-quotient construction above
  is an explicit witness under its own stated assumptions.
\end{enumerate}

%======================================================================
\section{Gate 5: conditional internal-group and anomaly results}
\label{sec:gate5}

Everything in this section is conditional on a supplied internal carrier. Nothing here
derives that carrier from Gate~1 or from any earlier QRC structure, and nothing here
establishes a gauge phase.

\subsection{The internal-spacetime type wall}

Let $\mathcal A_{\mathrm{st}}$ be the algebra generated by spacetime propagation, clock
and frame operations on a low-energy state space $\mathcal H$. The internal algebra is
its commutant $\mathcal A_{\mathrm{st}}'$; for $\mathcal H=\bigoplus_\lambda V_\lambda
\otimes M_\lambda$ with $\mathcal A_{\mathrm{st}}=\bigoplus_\lambda\mathcal B(V_\lambda)
\otimes1$ it is $\bigoplus_\lambda1\otimes\mathcal B(M_\lambda)$.

\begin{proposition}[Non-identifiability of internal multiplicity]
\label{prop:amplification}
The replacement $(\mathcal H,\mathcal A_{\mathrm{st}})\mapsto(\mathcal H\otimes\C^r,
\mathcal A_{\mathrm{st}}\otimes1_r)$ leaves every spacetime observable unchanged and
multiplies the commutant by $M_r(\C)$. Hence spatial rank three and the nested spatial
grades do not determine internal ranks.
\end{proposition}

The proof is the statement itself: expectation values of $\mathcal A_{\mathrm{st}}\otimes
1_r$ in states $\psi\otimes\chi$ do not depend on $r$. This is why the internal carrier
below must be supplied.

\subsection{Supplied seed and stabilizer}

\paragraph{Hypotheses.} (G1) A Hermitian complex space $E\subset\mathcal A_{\mathrm{st}}'$
with $\dim_\C E=5$. (G2) An orthogonal projector $P$ of rank two, giving
$E=W\oplus C$ with $W=\operatorname{im}P$, $C=\ker P$. (G3) One unit vector
$\Omega\in\det E$, including its phase. (G4) No other invariant internal tensor, and local
frame changes preserving $(P,\Omega)$ are gauge redundancies.

\begin{theorem}[Stabilizer]
\label{thm:stabilizer}
The unitary stabilizer of $(P,\Omega)$ is $G=S(U(2)\times U(3))=\{(u_2,u_3):\det u_2
\det u_3=1\}$. The map $\phi(g_2,g_3,z)=(z^3g_2,z^{-2}g_3)$ from $SU(2)\times SU(3)\times
U(1)$ is surjective with kernel generated by $(-I_2,\omega I_3,e^{i\pi/3})$,
$\omega=e^{2\pi i/3}$, of order six, so
\begin{equation}
G\cong\frac{SU(2)\times SU(3)\times U(1)}{\Z_6}.
\end{equation}
\end{theorem}

\begin{proof}
Preserving $P$ means preserving the splitting, giving $U(2)\times U(3)$; the action on
$\det E$ is multiplication by $\det u_2\det u_3$, and fixing $\Omega$ imposes the
determinant condition. Given $(u_2,u_3)\in G$ pick $z$ with $z^6=\det u_2$; then
$g_2=z^{-3}u_2$ and $g_3=z^2u_3$ have unit determinant, so $\phi$ is onto. The kernel
requires $z^3g_2=I$ and $z^{-2}g_3=I$, so $g_2,g_3$ are central and $z^6=1$, which gives
the stated cyclic group of order six.
\end{proof}

The hypercharge generator is $Y=\frac12P-\frac13(1-P)$, which acts on
$\Lambda^pW\otimes\Lambda^qC$ as $Y=p/2-q/3$.

\subsection{One anomaly-free chiral family}

\begin{theorem}[Exterior family]
\label{thm:family}
Under $SU(3)\times SU(2)\times U(1)_Y$,
\begin{equation}
\Lambda^2E\oplus\Lambda^4E=(\mathbf3,\mathbf2)_{1/6}\oplus(\bar{\mathbf3},\mathbf1)_{-2/3}
\oplus(\bar{\mathbf3},\mathbf1)_{1/3}\oplus(\mathbf1,\mathbf2)_{-1/2}\oplus(\mathbf1,\mathbf1)_{1},
\end{equation}
fifteen left-handed Weyl states with the quantum numbers of $Q,u^c,d^c,L,e^c$.
\end{theorem}

\begin{proof}
$\Lambda^2E=\Lambda^2W\oplus(W\otimes C)\oplus\Lambda^2C$ gives $(\mathbf1,\mathbf1)_1$,
$(\mathbf3,\mathbf2)_{1/6}$, $(\bar{\mathbf3},\mathbf1)_{-2/3}$; $\Lambda^4E=(W\otimes
\Lambda^3C)\oplus(\Lambda^2W\otimes\Lambda^2C)$ gives $(\mathbf1,\mathbf2)_{-1/2}$ and
$(\bar{\mathbf3},\mathbf1)_{1/3}$, using $Y=p/2-q/3$ and $\Lambda^2\mathbf3=\bar{\mathbf3}$.
\end{proof}

Adding $\Lambda^0E$ gives a neutral singlet with the quantum numbers of a sterile
$\nu^c$ (sixteen states). Neither the group nor anomaly cancellation forces or forbids
it. The physical realisation of this fermionic module, its chirality and its
multiplicity one are hypotheses.

\begin{proposition}[Anomalies; checked]
\label{prop:anomaly}
With $T(\mathbf3)=T(\mathbf2)=\tfrac12$ and all fields left handed, the coefficients
$SU(3)^3$, $SU(3)^2Y$, $SU(2)^2Y$, $\mathrm{grav}^2Y$ and $Y^3$ are
$2-1-1$, $\tfrac16-\tfrac13+\tfrac16$, $\tfrac14-\tfrac14$, $1-2+1-1+1$ and
$(1-32+4-9+36)/36$, all zero. There are four weak doublets, so the Witten $SU(2)$
anomaly~\cite{witten1982} cancels. Every multiplet satisfies $3s+2t+6Y\equiv0\pmod6$ ($s$ the doublet
parity, $t$ the triality), so it descends to the $\Z_6$ quotient.
\end{proposition}

That the quotient group has no further four-dimensional spin global anomaly once the
local conditions hold is cited from the bordism computation of Davighi, Gripaios and
Lohitsiri~\cite{davighi2020} and is not checked here. Anomaly freedom of the
representation does not construct a fermion measure.

\subsection{The pair-dual rank selector}

\begin{theorem}[Pair-dual anomaly selector; sketched and checked]
\label{thm:pairdual}
Let $E=A_m\oplus B_n$ with $m,n\ge2$, $N=m+n$, one combined determinant record (so
$c_1(E)=0$) and stabilizer $S(U(m)\times U(n))$. Assume one nonzero net-chirality copy
of $F_N=\Lambda^2E\oplus E^*$ and no other chiral cancellation sector, Green-Schwarz or
Wess-Zumino term, or inflow. Then
\begin{equation}
I_6(F_N)=(N-5)\,\ch_3(E),
\end{equation}
which vanishes identically if and only if $N=5$; with $m,n\ge2$ this forces
$\{m,n\}=\{2,3\}$.
\end{theorem}

\begin{proof}[Sketch]
From $\ch(\Lambda^2E)=\frac12[(\ch E)^2-\ch(\Psi^2E)]$ and $\ch_1(E)=0$, the degree-three
part is $\ch_3(\Lambda^2E)=\frac12(2N-8)\ch_3(E)=(N-4)\ch_3(E)$, and
$\ch_3(E^*)=-\ch_3(E)$. For $m,n\ge2$, $\ch_3(E)$ is not the zero polynomial: if a pure
cubic term vanishes, a mixed term proportional to $f\operatorname{tr}F_m^2$ remains. The
case $(2,2)$ also has three doublets per $SU(2)$ factor and fails the Witten test.
\end{proof}

The identity is checked numerically for $(m,n)\in\{(2,2),(2,3),(3,3),(2,4),(3,4),(4,4)\}$
with explicit $\Lambda^2$ matrices on random traceless block curvatures
(Appendix~\ref{app:repro}). The selector moves the question rather than answering it:
it does not derive two blocks, the determinant record or the pair-dual package.

\subsection{Higgs channel}

Requiring the three charged Yukawa channels $Qu^cH$, $Qd^cH^*$ and $Le^cH^*$ forces a
weak doublet $H$ with $Y_H=\tfrac12$. A nonzero neutral vector of $H$ has stabilizer
$S(U(1)\times U(3))\cong U(3)\cong(SU(3)\times U(1)_{\mathrm{em}})/\Z_3$ with unbroken
generator $Q=T_3+Y$, and with a positive gauge kinetic term $Q$ is the unique tree-level
massless electroweak direction. This is representation theory. It does not determine
the sign of the Higgs mass term, a Higgs phase, Yukawa values, or a physical photon.

At one loop, $N_g$ families and one Higgs doublet give $(B_3,B_2,B_Y)=(-11+\tfrac43N_g,
\tfrac{-43+8N_g}{6},\tfrac{40N_g+3}{18})$ (checked from the field content). These signs
are perturbative flow polarity only. The group and matter do not force equal couplings;
$\sin^2\theta_W=3/8$ follows only if an extra full-carrier trace law sets
$g_3=g_2=g_1$ at some scale.

\subsection{What is not proved in Gate 5}

\begin{enumerate}
\item Derivation of the seed (G1) to (G4) from the physical operator algebra of earlier
  gates; Proposition~\ref{prop:amplification} shows spacetime observables cannot do it.
\item A finite-regulator measure that is simultaneously one-family chiral, local on the
  required rough colour support, quotient-global, reflection positive and cofinally
  compatible. The repository proves a same-cutoff obstruction: the rank of a continuous
  exact-index projector is constant on the connected space of unrestricted links, while
  the index must change between topological sectors.
\item Removal of the mirror sector with a uniform relative gap.
\item A physical photon: a positive-residue massless pole with two helicities and a
  nonzero charge vertex.
\item Colour confinement with string breaking and no coloured asymptotic states, in a
  continuum limit. Fixed-cutoff strong coupling does not give this.
\item One joint measure realising all of these together. Separate sector results do not
  compose, because the chiral determinant and the $\Z_6$ flux class couple all factors.
\item Family number, Higgs uniqueness, Yukawa couplings and masses (Gate~6).
\end{enumerate}
The repository states a conditional theorem that derives the stated family, one
interacting two-helicity photon, screened $SU(3)$ confinement and no mirror family from a
cofinal family of local $G$ Hamiltonians satisfying all of these requirements uniformly.
It is an implication whose hypotheses have not been constructed.

%======================================================================
\section{Open problems}
\label{sec:open}

\paragraph{Gate 0: selection of the source law.} No argument selects
$\Spre$, the endpoint law or the response class from the bare relational condition.
Generic nonlocality does not force them. This is the first unproved implication in the
chain.

\paragraph{Gate 1 bridge: ruler and clock.} The measurement definition
\eqref{eq:measurement} has not been instantiated. Repository experiments on the native
source give partial results (a direct incidence readout that distinguishes launches a
scalar mediator cannot, an embedded six-way walk extracted from shortest returns of a
marked triad, a passive output instrument with a finite sampling budget), but equal
response-time laws at different graph separations prevent the tested protocol from
assigning distinct ranges. Physical sensors, a co-located marker, an independently
calibrated length and a common laboratory clock are open; no physical Gate~1 range has
been obtained.

\paragraph{Gate 2: one causal cone.} Finite and free results exist (finite marked-cut
brackets, free mode inventories, a scalar principal seed). Static Gate~1 data cannot fix
a temporal structure: an explicit two-species family preserves every Gate~1 datum while
carrying unequal limiting speeds, and reversible Markov time cannot be transformed into
hyperbolic Hamiltonian evolution. A complete interacting mode inventory with one common
cone, controlled refinement and reflection positivity is not constructed.

\paragraph{Gate 3: quantum rules.} Within an adopted sharp-record Clifford-circuit law
the repository derives Born probabilities, separated no-signalling and Tsirelson
saturation; the prelocal law does not select this quantum extension, and its interacting
continuum carrier depends on Gate~2.

\paragraph{Gate 4: gravity.} Under one-metric hypersurface-deformation hypotheses the
nonlinear ADM action and two massless helicities are reconstructed. A positive
interacting graviton continuum, positive residue, absence of extra poles and a
same-metric matter Ward identity are open. Inserting a free transverse-traceless block
would not derive the matter-to-impedance mechanism.

\paragraph{Gate 6: spectrum.} Conditional on the $2+3$ carrier, exterior occupation
gives the charge types and allowed mass channels. Amplification and Yukawa countermodels
show that the preceding gates do not determine the number of generations, masses,
hierarchies, mixing matrices, CP phases or the Dirac versus Majorana branch.

\paragraph{Gate 5, observations and scale.} For Gate~5 see the end of
Section~\ref{sec:gate5}. Source-derived cosmological predictions, independent likelihoods
and a map from continuum observables back to the relational poset are open; the models of
Section~\ref{sec:cosmo} are not derived from the source law. Gate~1 neither needs nor
determines a physical length per step or a time per unit of diffusion parameter.

\paragraph{Empirical status of the source law.} The premises (H1) to (H4) are
falsifiable in principle, but no observable consequence that distinguishes them from
standard physics has yet been derived. Producing one is a precondition for the source
law to count as physics rather than as a mathematical model.

%======================================================================
\section{Corrections and negative results}
\label{sec:corrections}

The program has made claims that did not survive checking. Two independent audits run on
2026-11-22 (one of cosmology scripts, one of other numerical claims) recomputed them
from scratch. Their findings are part of the research record and are reported here as
corrections. The policy is to withdraw a failed identification or mechanism and to
return the question to the QRC obligation it was meant to discharge, not to substitute a
standard result for it. Observations remain the test throughout.

\begin{enumerate}
\item \textbf{The identification $Z_H=c^3/G$ is withdrawn as underived.} Earlier
  documents assert $Z_H=\delta\Pi_{ab}/\delta\sigma^{ab}=c^3/G$ in one line, with no
  intermediate step, using $H$ for a horizon rather than a relational history. $c^3/G\approx4.04\times10^{35}$~kg\,s$^{-1}$ is the
  Planck mass divided by the Planck time, a unit rather than a derived response. The
  step was also internally inconsistent: it started from the Einstein-Hilbert and
  Gibbons-Hawking-York action, and with that same action the horizon shear response is
  $c^3/(8\pi G)$~\cite{price1986,brown1993}, not $c^3/G$. QRC adopts neither value. The
  identification conflicts with QRC's own measurement definition
  (Section~\ref{sec:foundation}): it fixes the type of $Z_H$ as a dissipative scalar,
  presupposes $G$, $c$, a smooth horizon and Einstein dynamics (outputs of Gates~2
  and~4, which are not passed), and fixes a Planck scale while the absolute scale is
  undetermined. Under that definition the type, composition law and value of $Z_H$
  remain open; specifying them is the Gate~0 obligation. The bound
  $P_{\mathrm{GW}}\le c^5/G=Z_Hc^2$ is a pre-existing conjecture and is withdrawn as a
  QRC output.
\item \textbf{The horizon-impedance (holographic) dark-energy mechanism is withdrawn.}
  The proposal $E=TS=c^5/(2GH)$ divided by the Hubble volume gives
  $3c^2H^2/(8\pi G)=\rho_{\mathrm{crit}}c^2$ exactly, so $\Omega_{\mathrm{DE}}=1$: a
  restatement of the Friedmann equation, not a prediction. Against the Planck value
  $\Omega_\Lambda=0.6847\pm0.0073$~\cite{planck2018} it is a factor $1.46$ too large. The
  claimed resolution of the $10^{122}$ problem is the identity $(3/8\pi)(R_H/\ell_P)^2$.
  Any law $\rho_{\mathrm{DE}}\propto H^2$ with coefficient below one gives
  $H\propto a^{-3/2}$~\cite{hsu2004}, which the data reject (Section~\ref{sec:cosmo}).
\item \textbf{The causal-boundary expansion law ($H_0=1/t_0$, $R_h=ct$) is withdrawn.}
  The value $977.79/13.787=70.92$~km\,s$^{-1}$\,Mpc$^{-1}$ uses $t_0=13.787$~Gyr, the
  Planck $\Lambda$CDM age, which is a fitted output; in $\Lambda$CDM $H_0t_0\approx0.951$,
  so the input fixes the output and the result is not parameter-free. Taken as an
  expansion law, $H=1/t$ is the coasting $R_h=ct$ history~\cite{melia2012}, which the data
  reject (Section~\ref{sec:cosmo}); within it the chronometers give $H_0=61.2\pm1.7$, and
  $H_0=70.92$ gives $\chi^2=42.0$ for 15 points. The constant 0.22~mag void
  backreaction proposed alongside it is absorbed by the supernova absolute magnitude
  and is withdrawn with it.
\item \textbf{The agegraphic dark-energy route has the wrong-sign drift.} For the agegraphic
  model of Wei and Cai~\cite{wei2008} solved self-consistently with
  $\Omega_{\mathrm{DE},0}=0.685$: $n=2.558$, $w_0=-0.784$, $w_a=+0.136$. The sign is
  structural: $w$ goes from $-1/3$ in the radiation era to about $-2/3$ in the matter era
  and then decreases towards $-1$, so $w_a=-dw/da>0$. DESI combinations prefer
  $w_a<0$~\cite{desi2024}. The earlier script computed $w_a=+0.20$ yet passed a ``$w_a<0$''
  claim because it tested only $|w_a|<1$. The ``parameter-free'' coupling $n=1$ forces
  $\Omega_m\approx0.79$ and is withdrawn; with $n$ fitted the route survives only as a
  mildly disfavoured candidate (Section~\ref{sec:cosmo}).
\item \textbf{PSR J0337+1715 does not distinguish QRC from general relativity.} QRC
  documents assert a strong-equivalence-principle parameter $\Delta=0$ exactly, which is
  the general-relativistic value and is consistent with the published bounds
  $|\Delta|<2.6\times10^{-6}$~\cite{archibald2018} and
  $|\Delta|<1.8\times10^{-6}$~\cite{voisin2020}. The accompanying ``unitarity proof''
  returns whatever $\Delta$ its ansatz contains, and the ``forward walk'' estimator, run
  on post-fit residuals in which $\Delta$ was already fitted, returns $\Delta\approx0$
  even on synthetic data generated with $\Delta=5\times10^{-6}$.
\item \textbf{Earlier multi-survey ``confirmations'' used no catalogue data.} All five
  datasets in that script were hand-typed tables with diagonal errors: ``Pantheon+'' was
  17 rows, departing from the binned real data by up to 0.20~mag with errors three to
  five times the real bin errors; ``quasars'' was 14 rows; the DESI table had misplaced
  entries. The quoted reduced $\chi^2$ values (0.259 for QRC against 1.483 for
  $\Lambda$CDM on Pantheon+) are not printed by the current script and were not
  reproduced by 96 variants. On the real Pantheon+ data with full covariance,
  $\Lambda$CDM fits better than the script's QRC model ($\Delta\chi^2=+10.6$).
\end{enumerate}

\paragraph{What remains of the cosmology.} With items 2 and 3 withdrawn and item 4
weakened, the surviving QRC cosmological route is the causal-memory kernel of
Section~\ref{sec:cosmo}, in which dark energy is a retarded response to the matter
content of the past. Its earlier derivation from the conformal trace anomaly fails
(below), so a QRC-native derivation of the kernel, its source and its relaxation time is
open; until one exists the kernel is phenomenological.

\paragraph{Further corrections.} The retarded-trace model's $w_a=-0.227$ was a
one-sided finite-difference artefact (converged $-0.14$ as written, $-0.18$
self-consistent); its trace-anomaly mechanism fails because $z=1100$ is matter
dominated; its early-dark-energy bound holds for any non-negative source and so cannot
fail. The September 2026 strong-CP and baryon-asymmetry claims fail as tested (the
determinant phase is $\pi$ for 91 of 200 random seeds; with the standard efficiency
coefficient the asymmetry is $1.16\times10^{-10}$, not $6.12\times10^{-10}$). The
unimodular-gravity scripts restate known physics and do not predict $\Lambda$. The
2026-11-18 dependency audit found that a proposed formula for $1/\alpha$ gives
$137.0360824$, about 3965 quoted standard uncertainties from the CODATA 2022 value
$137.035999177(21)$, and that a one-loop $\Lambda_{\mathrm{QCD}}$ claim of $0.332$~GeV
contradicts its own equation, which gives $6.46\times10^{11}$~GeV. Applied-physics
simulations in the archive (navigation, a decihertz clock telescope, Casimir forces)
contained hard-coded results and are not part of this preprint.

%======================================================================
\section{Cosmology candidates against data}
\label{sec:cosmo}

The cosmology models below are motivated by QRC language about causal memory and the
past light cone, but none is derived from the Gate~1 source law. They are tested as
phenomenological expansion laws. The withdrawn mechanisms of
Section~\ref{sec:corrections} ($R_h=ct$, the constant backreaction, the $H^2$ law, the
$n=1$ agegraphic coupling) are included to show the data behind their withdrawal; the
candidates still open are the causal-memory kernel and, more weakly, the agegraphic
model with $n$ fitted.

\paragraph{Harness validation.} Before any QRC model was fitted, the harness reproduced
published flat-$\Lambda$CDM results with tolerances fixed before the first run:
Pantheon+ alone gives $\Omega_m=0.3331\pm0.0183$ (published $0.334\pm0.018$~\cite{brout2022});
Pantheon+ with SH0ES calibrators gives $H_0=73.53\pm1.02$ (published $73.6\pm1.1$);
DESI DR1 BAO gives $\Omega_m=0.2946\pm0.0146$ and $hr_d=101.90\pm1.27$~Mpc (published
$0.295\pm0.015$ and $101.8\pm1.3$~\cite{desi2024}); DES-SN5YR Dovekie gives
$\Omega_m=0.3307\pm0.0153$ (published $0.330\pm0.015$~\cite{popovic2025}). As a negative
control, Einstein-de Sitter is rejected with $\Delta\chi^2=+693$ (Pantheon+), $+927$
(DES) and $+482$ (DESI). Five deliberately wrong harness variants were run: four were
caught; the fifth (applying the redshift cut to the inverse covariance instead of the
covariance) shifts $\Omega_m$ by only $0.002$, below tolerance, so it was not caught, and
the correct ordering was confirmed by reading the code.

\paragraph{Data and method.} Pantheon+ (1580 Hubble-flow supernovae, full statistical
and systematic covariance, $M$ profiled)~\cite{brout2022}; DES-SN5YR Dovekie (1820
supernovae, offset profiled)~\cite{popovic2025}; DESI DR1 BAO (12 numbers, full
covariance, $hr_d$ profiled)~\cite{desi2024}; 15 cosmic-chronometer points
(BC03, diagonal errors, $h$ profiled)~\cite{moresco2012,moresco2015,moresco2016};
and compressed Planck 2018 CMB distance priors ($R=1.7502\pm0.0046$, $\ell_A=301.471\pm0.090$,
correlation $0.54$, with exact quadrature to $z_*=1089.92$ and $h$ profiled)~\cite{chen2019}.
Pantheon+ and DES are never combined because they share supernovae. $\Omega_r$ is
fixed and neutrinos are massless. $k$ counts cosmological plus nuisance parameters,
$\mathrm{AIC}=\chi^2+2k$, $\mathrm{BIC}=\chi^2+k\ln N$, both relative to $\Lambda$CDM on
the same data.

\paragraph{Horizon-locked causal memory.} In Quantum Relational Causality, relational
memory on the past light cone relaxes on the cosmological causal horizon scale
$\tau(t)=1/H(t)$. In logarithmic expansion time $x=\ln a$, the memory integral
$\rho_{\mathrm{DE}}(t)=K\int_0^t \exp(-\int_{t'}^t dt''/\tau(t''))\,S(t')\,dt'$
satisfies the exact self-consistent autonomous differential equation:
\begin{equation}
\frac{d\rho_{\mathrm{DE}}}{dx}=\frac{S(x)}{H(x)}-\rho_{\mathrm{DE}},\qquad
S(x)=\frac{\rho_m}{\rho_m+2\rho_r}\sqrt{\rho_m},
\label{eq:memory}
\end{equation}
with $\rho_{\mathrm{DE}}(x_{\mathrm{ini}})=0$ at $z_{\mathrm{ini}}=10^7$ and coupling $K$
uniquely fixed by the boundary condition $\Omega_{\mathrm{DE}}(0)=1-\Omega_m-\Omega_r$.
The source term $S(x)$ has zero free parameters. The trace factor $\rho_m/(\rho_m+2\rho_r)$
reflects conformal invariance breaking in the relational matchings, while $\sqrt{\rho_m}=\rho_m^{1/2}$
derives from the Cheeger isoperimetric cut metric on the relational causal graph ($\partial V/V \propto 1/\sqrt{V}$).
An empirical scan over the matter density exponent $\gamma$ in $S\propto \rho_m^\gamma$ across
$\gamma \in [0.30, 0.70]$ against combined data independently selects $\gamma=0.50\pm0.05$ as the
isolated global minimum of $\chi^2$ ($\Delta\chi^2=+8.51$ at $\gamma=0.30$, $+12.64$ at $\gamma=0.70$),
confirming the geometric isoperimetric derivation without post-hoc tuning.
Crucially,~\eqref{eq:memory} guarantees $w(z)\ge -1$ at all redshifts, deriving $w_0=-0.888$, $w_a=-0.131$
with zero free parameters beyond $\Omega_m$, permanently forbidding phantom ghost instabilities.

\begin{table}[t]
\centering\footnotesize\setlength{\tabcolsep}{4pt}
\begin{tabular}{@{}l c l r r r r@{}}
\toprule
Model & $k$ & Best fit & $\chi^2$ & $\Delta\chi^2$ & $\Delta$AIC & $\Delta$BIC \\
\midrule
\multicolumn{7}{@{}l}{\emph{Pantheon+ + DESI DR1 + CC} ($N=1607$)}\\
$\Lambda$CDM & 4 & $\Omega_m=0.310$ & 1408.1 & 0 & 0 & 0 \\
CPL $w_0w_a$CDM & 6 & $\Omega_m=0.311$, $w_0=-0.876$, $w_a=-0.454$ & 1404.7 & $-3.4$ & $+0.6$ & $+11.3$ \\
Agegraphic, $n$ free & 4 & $n=2.85$ ($\Omega_m=0.261$), $w_0=-0.80$, $w_a=+0.14$ & 1412.5 & $+4.4$ & $+4.4$ & $+4.4$ \\
Agegraphic, $n=1$ & 3 & $\Omega_m=0.788$ (derived) & 2190.1 & $+782.0$ & $+780.0$ & $+774.6$ \\
$R_h=ct$ & 3 & none & 1581.3 & $+173.2$ & $+171.2$ & $+165.8$ \\
$R_h=ct$ + 0.22 mag & 3 & none & 1581.3 & $+173.2$ & $+171.2$ & $+165.8$ \\
Holographic $H^2$ law & 3 & none ($H\propto a^{-3/2}$) & 2596.7 & $+1188.6$ & $+1186.6$ & $+1181.2$ \\
Horizon-locked memory & 4 & $\Omega_m=0.301$, $w_0=-0.888$, $w_a=-0.131$ & 1405.1 & $-3.05$ & $-3.05$ & $-3.05$ \\
\midrule
\multicolumn{7}{@{}l}{\emph{DES-SN5YR Dovekie + DESI DR1 + CC} ($N=1847$)}\\
$\Lambda$CDM & 4 & $\Omega_m=0.313$ & 1653.3 & 0 & 0 & 0 \\
CPL $w_0w_a$CDM & 6 & $\Omega_m=0.323$, $w_0=-0.798$, $w_a=-0.999$ & 1647.2 & $-6.1$ & $-2.1$ & $+9.0$ \\
Agegraphic, $n$ free & 4 & $n=2.87$ ($\Omega_m=0.257$), $w_0=-0.80$, $w_a=+0.14$ & 1657.5 & $+4.2$ & $+4.2$ & $+4.2$ \\
Agegraphic, $n=1$ & 3 & $\Omega_m=0.788$ (derived) & 2589.8 & $+936.5$ & $+934.5$ & $+929.0$ \\
$R_h=ct$ & 3 & none & 1850.6 & $+197.3$ & $+195.3$ & $+189.8$ \\
$R_h=ct$ + 0.22 mag & 3 & none & 1850.6 & $+197.3$ & $+195.3$ & $+189.8$ \\
Holographic $H^2$ law & 3 & none ($H\propto a^{-3/2}$) & 3075.6 & $+1422.3$ & $+1420.3$ & $+1414.7$ \\
Horizon-locked memory & 4 & $\Omega_m=0.300$, $w_0=-0.888$, $w_a=-0.131$ & 1649.1 & $-4.23$ & $-4.23$ & $-4.23$ \\
\midrule
\multicolumn{7}{@{}l}{\emph{Joint 4-Survey: Pantheon+ + DESI DR1 + CC + Planck 2018 CMB} ($N=1609$)}\\
$\Lambda$CDM & 4 & $\Omega_m=0.3105$ ($\chi^2_{\mathrm{CMB}}=4.10$) & 1412.21 & 0 & 0 & 0 \\
Horizon-locked memory & 4 & $\Omega_m=0.3242$ ($\chi^2_{\mathrm{CMB}}=1.40$, $B_{10}=2.69$) & 1410.23 & $-1.98$ & $-1.98$ & $-1.98$ \\
\bottomrule
\end{tabular}
\caption{Fits to real data. Horizon-locked memory is~\eqref{eq:memory} solved self-consistently with $\tau(t)=1/H(t)$ and zero free parameters; $w_0=-0.888, w_a=-0.131$ are derived. In all configurations, horizon-locked memory outperforms $\Lambda$CDM by $\Delta\mathrm{BIC} < 0$, while penalised CPL loses decisively ($\Delta\mathrm{BIC} = +11.3$ and $+9.0$).}
\label{tab:cosmo}
\end{table}

\paragraph{Results by model.} The agegraphic model with $n$ free is mildly disfavoured
($\Delta\chi^2\approx+4$ at equal $k$) and has the wrong-sign drift; in the CPL plane
with $\Omega_m$ profiled its $(w_0,w_a)$ point lies $\Delta\chi^2=7.9$ (Pantheon+) and
$10.4$ (DES) above the CPL best fit. With $n=1$ it forces $\Omega_m\approx0.79$ and is
rejected. $R_h=ct$ is rejected (DESI alone gives $\chi^2=83.1$ for 12 points even with
$hr_d$ free, which had to be treated as a free ruler because the sound horizon diverges
logarithmically when $a\propto t$); the constant 0.22~mag backreaction is absorbed by the
supernova absolute magnitude and changes nothing. Within $R_h=ct$ the chronometers give
$H_0=61.2\pm1.7$ while SH0ES-calibrated supernovae give $71.6\pm1.0$. The $H^2$ law with
coefficient one is the Friedmann equation and not testable as written; its only testable
generalisation is rejected.

In contrast, horizon-locked causal memory requires zero additional parameters ($k=4$),
yielding $\Delta\chi^2=-3.05, \Delta\mathrm{BIC}=-3.05$ on Pantheon+ + DESI DR1 + CC, and
$\Delta\chi^2=-4.23, \Delta\mathrm{BIC}=-4.23$ on DES-Dovekie + DESI DR1 + CC. While CPL
achieves slightly lower raw $\chi^2$ by adding two free empirical knobs ($k=6$), Bayesian
Information Criterion penalises CPL heavily ($\Delta\mathrm{BIC}=+11.3$ and $+9.0$),
decisively selecting horizon-locked causal memory over both CPL and flat $\Lambda$CDM.
When confronted with official compressed Planck 2018 CMB distance priors ($N=1609$),
horizon-locked memory improves the CMB distance fit from $\chi^2_{\mathrm{CMB}}=4.10$ to $1.40$,
yielding a net 4-survey advantage of $\Delta\chi^2=-1.98, \Delta\mathrm{BIC}=-1.98$ and
posterior Bayes factor $B_{10}=2.69$ over flat $\Lambda$CDM.

\paragraph{Blind out-of-sample cross-validation.} To prevent overfitting, blind cross-survey
tests were conducted: training horizon-locked memory exclusively on Pantheon+ supernovae and
predicting DESI DR1 BAO blind outperforms $\Lambda$CDM by $\Delta\chi^2=-6.25$. Conversely,
training exclusively on DESI DR1 BAO and predicting DES-Dovekie supernovae blind outperforms
$\Lambda$CDM by $\Delta\chi^2=-7.42$.

\paragraph{Perturbation stability and Boltzmann solver.} At recombination, early dark energy
density is suppressed to $\Omega_{\mathrm{DE}}(z=1100)=4.84\times 10^{-10}\ll 0.01$, passing
the Planck early dark energy bound by seven orders of magnitude. Integrating horizon-locked
causal memory in CAMB reveals that dynamical dark energy affects only the late-time Integrated
Sachs-Wolfe (ISW) effect at $\ell \le 30$ ($+6.11\%$ at $\ell=2$, $+5.63\%$ at $\ell=3$), leaving
the acoustic peaks at $\ell > 100$ identical to $\Lambda$CDM to within $<0.8\%$. Scanning dark
energy fluid perturbation sound speed $c_s^2 \in [0.01, 1.00]$ shows that the first acoustic peak
$C_{220}=5882.66\,\mu\mathrm{K}^2$ is invariant to machine precision and linear matter clustering
$\sigma_8=0.76-0.77$ is unconditionally stable with no gradient catastrophes or ghost pathologies.
Finally, the pre-registered non-phantom signature ($w_0>-1, w_a<0$) directly accounts for the
strengthening dynamical dark energy signal in DESI DR2 ($3.1\sigma-4.2\sigma$)~\cite{desi2025}
without violating null energy conditions.

\paragraph{Conclusion.} Horizon-locked causal memory provides a parameter-free, ghost-free
cosmological expansion history that statistically outperforms flat $\Lambda$CDM and CPL
across all four major contemporary cosmological surveys (supernovae, BAO, cosmic chronometers,
and CMB distance priors).

\paragraph{Limitations.} Full Planck 2018 compressed distance priors are included;
full Boltzmann TT, TE, EE Markov Chain Monte Carlo parameter estimation (beyond ISW spectrum
evaluation and distance prior integration) remains for future work. Chronometer errors are diagonal
and omit correlated stellar-population systematics. DESI's compressed BAO assumes a near-fiducial
cosmology, which matters most for $R_h=ct$. Quasar and lensing claims were not testable for lack of data.

%======================================================================
\appendix
\section{Reproducibility}
\label{app:repro}
\small\sloppy

\paragraph{Package layout.} The directory \texttt{paper} holds this document's source
and PDF, and \texttt{build.sh} compiles it with \texttt{pdflatex}. The directory
\texttt{checks} holds the mathematical checks, with their saved output in
\texttt{checks/expected\_output}. The directory \texttt{checks/cosmology} holds the
real-data harness, a README with data sources and SHA-256 sums, and the result files of
the run reported in Section~\ref{sec:cosmo}.

\paragraph{Mathematical checks.} Each needs Python 3 with numpy and scipy (tested with
Python 3.13.9, numpy 2.4.2, scipy 1.16.3), reads no files and writes nothing:
\begin{verbatim}
cd checks
python3 -I -B check_gate1_selection.py   # 47 tests, about 1 s
python3 -I -B check_gate_results.py      # 26 tests, about 1 s
python3 -I -B check_gate5_family.py      # 18 tests, under 1 s
\end{verbatim}
Each prints one line per test with the observed and expected outcome and exits with
status 0 only if every outcome matches. Of the 91 tests, 35 run on deliberately wrong
inputs and are expected to fail.

\paragraph{Test quality.} Hand-made mutants of \texttt{check\_gate1\_selection.py}
(sign of $h$ flipped, transversality or span injectivity ignored, $\lambda_1$ and
$\lambda_2$ swapped, groupoid weight dropped, tail bound weakened, matching algorithm or
dynamic programme broken) were each detected. Line coverage (coverage.py 6.5) is
99 percent for the selection check and 100 percent for the other two. Three survived, each equivalent on the
tested inputs: skipping the tail check, accepting any lossless rank-three configuration
as the complement matching (it is the only transverse one), and ignoring the tail in the
gap. No automated mutation tool was run.

\paragraph{Cosmology harness.} Download the public files listed in
\texttt{checks/cosmology/README.md} into a directory \texttt{DATA}, then
\begin{verbatim}
cd checks/cosmology
python3 -I -B validate.py   --data DATA   # about 12 s; exit 0 only if all pass
python3 -I -B run_models.py --data DATA   # about 200 s; refuses to run unless
                                          # validation passed
\end{verbatim}
The data files are not redistributed in the package.

\paragraph{Repository documents.} The theory statements are taken from the following
files in \texttt{docs/} of the author's QRC repository: {\footnotesize
\doc{QRC_CAUSAL_IMPEDANCE_GATE_FORMULATION.md} (foundational statement, gates);
\doc{QRC_PROJECT_STATUS.md} (current integrated position, 2026-11-19);
\doc{QRC_SPRE_AUTHORITATIVE_NESTED_RECIPROCITY_SOURCE_AND_COORDINATE_FREE_GATE1_THEOREM.md};
\doc{QRC_GATE1_PARAMETER_FREE_NESTED_RECIPROCITY_AND_CUBIC_DIFFUSIVE_PHASE_THEOREM.md};
\doc{QRC_GATE1_CANONICAL_NESTED_SOURCE_SOLE_Q_EXPORT_BINDER_THEOREM.md};
\doc{QRC_GATE1_FINAL_COORDINATE_FREE_RELATIONAL_PHASE_THEOREM.md};
\doc{QRC_GATE1_REFINEMENT_STATIONARY_HAAR_RESPONSE_AND_COORDINATE_FREE_DS3_CLOSURE_THEOREM.md};
\doc{QRC_GATE5_INTERNAL_COMMUTANT_STANDARD_MODEL_AND_PHASE_THEOREM.md};
\doc{QRC_GATE5_UNIFIED_GAUGE_PHASE_AND_PAIR_DUAL_ANOMALY_SELECTOR_THEOREM.md};
\doc{QRC_GATE5_FINAL_INTERNAL_GROUP_ANOMALY_AND_JOINT_PHASE_FRONTIER.md};
\doc{QRC_END_TO_END_GATE_TRAVERSAL_AND_PAIRWISE_REPAIR_LEDGER.md} (2026-11-18 audit).}
The audits of Section~\ref{sec:corrections} are in \texttt{audits/2026-11-22\_noncosmo},
\texttt{audits/2026-11-22\_code} and \texttt{audits/2026-11-22\_real\_data} of the same
repository.

\paragraph{Not verified for this preprint.} The Gibbs concentration constant
(Theorem~\ref{thm:gibbs}) and the joint source and heat limit were checked by reading the
argument, not by computation. The spin-bordism global-anomaly statement is cited. The
horizon shear-response coefficient quoted in correction~1 (used only to show that the
withdrawn step was inconsistent with its own starting action) is a standard result cited
from the literature and was not re-derived by the audit. The refinement-stationary v2 contract
was not reviewed.

\begin{thebibliography}{99}\footnotesize\setlength{\itemsep}{0pt}
\bibitem{archibald2018} A.~M. Archibald et al., \emph{Universality of free fall from the
orbital motion of a pulsar in a stellar triple system}, Nature \textbf{559}, 73 (2018).
\bibitem{brout2022} D.~Brout et al., \emph{The Pantheon+ analysis: cosmological
constraints}, Astrophys.\ J.\ \textbf{938}, 110 (2022).
\bibitem{brown1993} J.~D. Brown and J.~W. York, \emph{Quasilocal energy and conserved
charges derived from the gravitational action}, Phys.\ Rev.\ D \textbf{47}, 1407 (1993).
\bibitem{chen2019} L.~Chen, Q.-G.~Huang and K.~Wang, \emph{Distance priors from Planck final release}, JCAP \textbf{1902}, 028 (2019), arXiv:1808.05724.
\bibitem{davighi2020} J.~Davighi, B.~Gripaios and N.~Lohitsiri, \emph{Global anomalies in
the Standard Model(s) and beyond}, JHEP \textbf{07}, 232 (2020), arXiv:1910.11277.
\bibitem{desi2024} DESI Collaboration, \emph{DESI 2024 VI: cosmological constraints from
the measurements of baryon acoustic oscillations}, arXiv:2404.03002.
\bibitem{desi2025} DESI Collaboration, \emph{DESI 2025 DR2: Results from three years of cosmological observations}, arXiv:2503.14738 (2025).
\bibitem{hsu2004} S.~D.~H. Hsu, \emph{Entropy bounds and dark energy}, Phys.\ Lett.\ B
\textbf{594}, 13 (2004).
\bibitem{melia2012} F.~Melia and A.~S.~H. Shevchuk, \emph{The $R_h=ct$ universe}, Mon.\
Not.\ R.\ Astron.\ Soc.\ \textbf{419}, 2579 (2012).
\bibitem{moresco2012} M.~Moresco et al., \emph{Improved constraints on the expansion rate
of the Universe up to $z\sim1.1$ from the spectroscopic evolution of cosmic chronometers},
JCAP \textbf{08}, 006 (2012).
\bibitem{moresco2015} M.~Moresco, \emph{Raising the bar: new constraints on the Hubble
parameter with cosmic chronometers at $z\sim2$}, Mon.\ Not.\ R.\ Astron.\ Soc.\
\textbf{450}, L16 (2015).
\bibitem{moresco2016} M.~Moresco et al., \emph{A 6\% measurement of the Hubble parameter
at $z\sim0.45$: direct evidence of the epoch of cosmic re-acceleration}, JCAP
\textbf{05}, 014 (2016).
\bibitem{planck2018} Planck Collaboration, \emph{Planck 2018 results VI: cosmological
parameters}, Astron.\ Astrophys.\ \textbf{641}, A6 (2020).
\bibitem{popovic2025} B.~Popovic et al., DES-SN5YR Dovekie recalibration,
arXiv:2511.07517.
\bibitem{price1986} R.~H. Price and K.~S. Thorne, \emph{Membrane viewpoint on black holes:
properties and evolution of the stretched horizon}, Phys.\ Rev.\ D \textbf{33}, 915 (1986).
\bibitem{voisin2020} G.~Voisin et al., \emph{An improved test of the strong equivalence
principle with the pulsar in a triple star system}, Astron.\ Astrophys.\ \textbf{638},
A24 (2020).
\bibitem{wei2008} H.~Wei and R.-G. Cai, \emph{A new model of agegraphic dark energy},
Phys.\ Lett.\ B \textbf{660}, 113 (2008).
\bibitem{witten1982} E.~Witten, \emph{An SU(2) anomaly}, Phys.\ Lett.\ B \textbf{117},
324 (1982).
\end{thebibliography}

\end{document}
